% Results-section fragment for the Goodwin oscillator benchmarks.
% Figure boxes are intentional placeholders. Replace each box with the final
% \includegraphics command when the manuscript figure files are frozen.

\section{Simulated case studies}
\label{sec:results-goodwin}

We assessed the method on the three-state Goodwin oscillator
\begin{align}
    \dot{x} &= \frac{1}{1+z^{9.5}}-0.2x, \\
    \dot{y} &= x-0.15y, \\
    \dot{z} &= y-0.1z.
    \label{eq:goodwin-ground-truth}
\end{align}
This system provides a compact but nontrivial identification problem: its
feedback is strongly nonlinear, its trajectories are oscillatory, and the
correct right-hand side is sparse relative to the candidate libraries used
below.  Six trajectories were integrated on $t\in[0,100]$ with sampling
interval $\Delta t=0.1$.  The initial conditions were fixed before running the
experiments and split into three training trajectories, one validation
trajectory, and two test trajectories.  The validation and test initial
conditions were therefore never used to fit the coefficients.

Independent Gaussian noise was added to each state along each trajectory.  A
nominal noise level $\eta$ denotes a noise standard deviation equal to
$\eta$ times the empirical standard deviation of the corresponding clean
state.  Concentrations were clipped at $10^{-10}$ after noise addition to
preserve the positive domain required by the nonlinear candidate functions.
Five deterministic noise and optimizer seeds, $\{11,22,33,44,55\}$, were used
for every reported condition.

For each outer candidate, sequentially thresholded least squares (STLSQ) fitted
the linear coefficient matrix.  The multi-objective search minimized the
largest state-wise training trajectory mean-squared error and the number of
coefficients exceeding $10^{-6}$ in magnitude.  NSGA-II used a population of
80 for 30 generations, corresponding to 2400 candidate evaluations per run.
The final model was selected from the Pareto set by its rollout error on the
held-out validation trajectory, with sparsity used to break ties.  The test
trajectories were evaluated only after this selection.  Since the data are
simulated, validation and test normalized root-mean-square errors (NRMSEs) were
computed against the corresponding clean trajectories.  The reported rollout
NRMSE is the mean, over the three states, of the state-wise RMSE normalized by
the clean state standard deviation.

\subsection{Noise robustness and weak-form optimization}
\label{sec:results-goodwin-noise}

The first experiment isolated the effect of the regression formulation and
weak-test-function optimization while keeping the nonlinear Hill exponent
fixed at its ground-truth value.  The overcomplete library was
\begin{equation}
\begin{split}
\Theta(X)=\{&x,y,z,x^2,y^2,z^2,xy,xz,yz,\\
&h(x),h(y),h(z)\},
\qquad h(s)=\frac{1}{1+s^{9.5}}.
\end{split}
\label{eq:figure2-library}
\end{equation}
Thus, 12 candidate functions generated 36 possible coefficients across the
three equations, of which only six were nonzero in the ground truth.  No
constant term was included.

Four identification strategies were compared at noise levels of 0, 5, 10, 15,
and 20\%:
\begin{enumerate}
    \item \emph{strong}: Savitzky--Golay (SG) state smoothing and derivative
    estimation with window length 21 and polynomial order 3;
    \item \emph{weak, fixed}: a polynomial weak test function with degree 2,
    window length 61, and stride 20;
    \item \emph{weak, optimized}: joint optimization of test-function degree,
    window length, and stride; and
    \item \emph{weak, optimized + SG}: the same optimized weak formulation
    applied after the fixed SG state preprocessing used by the strong method.
\end{enumerate}
All four strategies optimized the STLSQ ridge parameter
$\alpha\in[10^{-6},1]$ and sparsification threshold
$\lambda\in[10^{-3},1]$.  The optimized weak variants additionally searched
degree $d\in[2,6]$, window length $L\in[21,151]$, and stride
$s\in[5,40]$, with integer-valued settings enforced by the formulation.  For
the SG-preprocessed weak variant, both regression and its training-trajectory
objective used the filtered training signal.  The SG window and polynomial
order themselves were not optimized in this experiment.

\begin{figure*}[t]
    \centering
    \fbox{\parbox[c][48mm][c]{0.94\textwidth}{\centering
    \textbf{Figure placeholder: noise robustness.}\\[1ex]
    Median and seed-level test rollout NRMSE, coefficient error, and
    false-positive count for the four formulations from 0 to 20\% noise.}}
    \caption{Goodwin noise-robustness experiment using the complete
    degree-two polynomial and fixed-Hill candidate library.  Points denote
    individual seeds and lines denote medians.}
    \label{fig:goodwin-noise-robustness}
\end{figure*}

Table~\ref{tab:goodwin-rollout-nrmse} reports the median held-out test NRMSE\@.
Without noise, both optimized weak variants reproduced the dynamics almost
exactly, whereas the strong formulation retained several unnecessary terms and
had a substantially larger rollout error.  As noise increased, SG
preprocessing produced a clear advantage for the optimized weak formulation.
At 20\% noise its median NRMSE was 0.1265, compared with 0.1986 for optimized
weak without SG, 0.2748 for the strong formulation, and 0.3140 for the fixed
weak formulation.  The SG-preprocessed optimized weak method also had the
smallest median coefficient error at every nonzero noise level.

\begin{table*}[t]
    \centering
    \caption{Median test rollout NRMSE across five seeds in the noise-robustness
    experiment.}
    \label{tab:goodwin-rollout-nrmse}
    \begin{tabular}{lccccc}
        \hline
        Method & 0\% & 5\% & 10\% & 15\% & 20\% \\
        \hline
        Strong & 0.1413 & 0.1402 & 0.1308 & 0.2070 & 0.2748 \\
        Weak, fixed & 0.00444 & 0.0781 & 0.1563 & 0.2347 & 0.3140 \\
        Weak, optimized & $8.94\times10^{-5}$ & 0.0654 & 0.1568 & 0.1844 & 0.1986 \\
        Weak, optimized + SG & $2.11\times10^{-4}$ & 0.00512 & 0.0274 & 0.0844 & 0.1265 \\
        \hline
    \end{tabular}
\end{table*}

The support results reinforce this dynamical comparison.  For optimized weak
with SG, all five seeds recovered the exact six-term support at 5, 10, and
15\% noise.  Four of five seeds recovered it at both 0 and 20\% noise.  At
0\% noise, one seed retained only negligible leakage: $8.0\times10^{-6}$ for
$yz$ and $-3.2\times10^{-7}$ for $z^2$ in the $\dot z$ equation.  At 20\%
noise, one seed selected a genuine spurious term, $-0.1167xy$ in the
$\dot x$ equation; this run also had the largest test NRMSE at that noise
level (0.2991).  Consequently, a median coefficient heat map alone would hide
this failure.  We therefore report both coefficient magnitude and selection
frequency across seeds.

\begin{figure*}[t]
    \centering
    \fbox{\parbox[c][48mm][c]{0.94\textwidth}{\centering
    \textbf{Figure placeholder: complete coefficient recovery.}\\[1ex]
    Ground-truth and recovered coefficient matrices for all 12 library terms,
    accompanied by seed-wise term-selection frequencies.}}
    \caption{Coefficient and support recovery for optimized weak identification
    with SG preprocessing.  The complete library is shown so that rejected
    squares, interactions, and incorrectly placed Hill terms remain visible.}
    \label{fig:goodwin-complete-coefficients}
\end{figure*}

Taken together, these results show that the benefit is not caused by restricting
the library to the known governing terms.  The algorithm was required to reject
30 zero coefficients, including all quadratic distractors and Hill functions
applied to the wrong states.  Fixed SG preprocessing substantially improved the
noise robustness of the optimized weak formulation, although the isolated
20\% failure shows that exact support recovery is not guaranteed for every
noise realization.

\subsection{Joint optimization of parametric library functions}
\label{sec:results-goodwin-parametric}

The second experiment tested whether the nonlinear library could be optimized
together with the weak formulation and sparse solver.  Based on the preceding
noise study, only the SG-preprocessed optimized weak method was used.  The
noise level was fixed at 10\%, and the library was enlarged to
\begin{equation}
\begin{split}
\Theta(X)=\{&x,y,z,x^2,y^2,z^2,xy,xz,yz,\\
&h_{\gamma_1}(x),h_{\gamma_1}(y),h_{\gamma_1}(z),\\
&e_{\gamma_2}(x),e_{\gamma_2}(y),e_{\gamma_2}(z)\},
\end{split}
\label{eq:figure3-library}
\end{equation}
where
\begin{equation}
    h_{\gamma_1}(s)=\frac{1}{1+s^{\gamma_1}},
    \qquad
    e_{\gamma_2}(s)=\exp(-\gamma_2 s).
\end{equation}
This library contains 15 functions and 45 possible coefficients.  The seven
outer optimization variables were the Hill exponent
$\gamma_1\in[4,12]$, exponential rate $\gamma_2\in[0.1,10]$, weak
degree, window length, and stride, and the STLSQ parameters $\lambda$ and
$\alpha$.  The population, generation count, train--validation--test split,
and SG preprocessing were unchanged from the first experiment.

An identifiability distinction is essential here.  The Goodwin ground truth
contains $h_{9.5}(z)$ but no exponential term.  Hence $\gamma_1$ has a
ground-truth recovery target only when a Hill coefficient is active, whereas
$\gamma_2$ has no ground-truth target at all.  If STLSQ removes all exponential
coefficients, changing $\gamma_2$ leaves the identified model unchanged; its
numerical value must therefore not be interpreted as a recovered physical
parameter.

\begin{figure*}[t]
    \centering
    \fbox{\parbox[c][52mm][c]{0.94\textwidth}{\centering
    \textbf{Figure placeholder: joint parametric-library experiment.}\\[1ex]
    Representative held-out dynamics; ground-truth and optimized Hill curves;
    selected exponential candidates; and the complete coefficient matrix.}}
    \caption{Joint optimization of the parametric library, weak formulation,
    and STLSQ solver using SG-preprocessed data at 10\% noise.  Faint nonlinear
    curves show seed-level results; the emphasized curve shows the conditional
    median over runs that selected the corresponding function family.}
    \label{fig:goodwin-parametric-library}
\end{figure*}

The correct Hill family was selected in three of five runs.  Conditional on
Hill selection, the recovered exponents were 9.4388, 10.0499, and 10.6703,
with median 10.0499 compared with the ground truth 9.5.  These three runs also
recovered the exact six-term support, with no polynomial or exponential false
positives.  Their test NRMSEs were 0.0568, 0.0290, and 0.0291, respectively.

In the remaining two runs, an exponential function of $z$ replaced the true
Hill feedback in the $\dot x$ equation.  The selected exponential rates were
$\gamma_2=6.0438$ and 5.6141.  One of these runs additionally retained
$z^2$ and $yz$ in the $\dot z$ equation.  Across all five seeds, the true
linear terms were selected with frequency 100\%, $h_{\gamma_1}(z)$ was
selected with frequency 60\%, and $e_{\gamma_2}(z)$ was selected with
frequency 40\%.  The median test NRMSE over all five runs was 0.0568, with
range 0.0290--0.1245.

\begin{figure}[t]
    \centering
    \fbox{\parbox[c][48mm][c]{0.90\columnwidth}{\centering
    \textbf{Figure placeholder: term-selection stability.}\\[1ex]
    Selection frequency of every library term in each state equation across
    the five seeds.}}
    \caption{Term-selection stability in the joint parametric-library
    experiment.  The competition between Hill and exponential feedback is
    visible in the $\dot x$ equation and is not apparent from median
    coefficients alone.}
    \label{fig:goodwin-parametric-support}
\end{figure}

This experiment demonstrates both the capability and the limitation of joint
library optimization.  When the correct functional family was selected, the
Hill exponent and sparse governing equations were recovered accurately from
noisy data.  However, over the finite state range represented by the training
trajectories, a rapidly decaying exponential could provide a competitive
surrogate for the Hill feedback.  Good held-out rollout accuracy therefore did
not always imply recovery of the correct nonlinear mechanism.  This distinction
motivates reporting family-selection frequency and conditional parameter error,
rather than only the best trajectory fit or an unconditional distribution of
the optimized library parameters.
